%% Elements of Nested Fibrational Cosmology
%% Book I: Primitive Relational Foundations
%% Final-Pass Canonical Draft
%%
%% Status legend:
%%   [U]  Unconditional  — proved from standing definitions and prior Unc results
%%   [C]  Conditional    — depends on explicitly declared hypotheses
%%   [D]  Definition-structural — structural, no proof obligation
%%   [R]  Remark only    — expository, non-load-bearing
%%   [O]  Open obligation — not yet proved

\documentclass[12pt,oneside]{amsbook}

\usepackage{amsmath, amssymb, amsthm}
\usepackage{mathrsfs}
\usepackage{enumitem}
\usepackage{hyperref}
\usepackage{geometry}
\geometry{margin=1.2in}

%% ---------------------------------------------------------------
%% Theorem environments
%% ---------------------------------------------------------------

\theoremstyle{plain}
\newtheorem{theorem}{Theorem}[section]
\newtheorem{lemma}[theorem]{Lemma}
\newtheorem{proposition}[theorem]{Proposition}
\newtheorem{corollary}[theorem]{Corollary}

\theoremstyle{definition}
\newtheorem{definition}[theorem]{Definition}
\newtheorem{standingrule}[theorem]{Standing Rule}
\newtheorem{commonnotion}[theorem]{Common Notion}

\theoremstyle{remark}
\newtheorem{remark}[theorem]{Remark}
\newtheorem{scholium}[theorem]{Scholium}

%% Status tag macro
\newcommand{\status}[1]{\textup{[#1]}}

%% ---------------------------------------------------------------
%% Notation
%% ---------------------------------------------------------------
%%   L              finite relational signature
%%   \Omega         configuration space
%%   \mathcal{T}_n  licensed test family at stage n
%%   \sim_n         observational equivalence at stage n
%%   \Sigma_n       observational quotient at stage n  (\Omega/{\sim_n})
%%   \Sigma_\infty  stabilized observational quotient (mnemonic only)
%%   C_n            class in \Sigma_n
%%   \delta(C_n)    one-step defect number
%%   \Delta         defect ledger
%%   S(\Omega)      structural entropy  = log|\Sigma_\infty|
%%   K_0            collar alphabet size (prime by SR~\ref{sr:alphabet-prime})
%%   \mathcal{E}    admissible extension class
%%   E_{\max}       maximal admissible class (Thm~\ref{thm:emax-unique})

\begin{document}

\title{Book I\\[6pt]Primitive Relational Foundations}
\author{Elements of Nested Fibrational Cosmology}
\date{}
\maketitle

\tableofcontents

\bigskip

%% ---------------------------------------------------------------
\section*{I.0\quad Purpose}
%% ---------------------------------------------------------------

Book~I fixes the primitive admissible arena of the canonical spine.
It begins from a purely finite relational setting and develops only
the minimal theorem-bearing grammar needed for admissible tests,
observational equivalence, quotient structure, refinement, finite
stabilization, and primitive defect bookkeeping.

No geometric, continuum, probabilistic, or branch-specific structure
is primitive here.  Such language, if it later appears in the spine,
must be licensed by later books.  Book~I therefore establishes the
first governing rule of the program:

\begin{quote}
\textit{What has not yet been derived is not yet licensed.}
\end{quote}

\noindent
Book~I does \emph{not} prove locality stability, source universality,
branch legitimacy, continuum correspondence, or any physical
interpretation.  It proves only the primitive theorem-bearing kernel:
admissible tests, observational quotient, finite refinement
stabilization, and defect bookkeeping.

%% ---------------------------------------------------------------
\section{Primitive Stance}
%% ---------------------------------------------------------------

We work in a purely finite relational setting throughout this book.
Observable content is generated only by admissible internal tests;
token-level identity has no canonical force beyond membership in a
finite site set.

\begin{standingrule}[\status{D}]\label{sr:no-primitive}
No primitive spacetime, metric, probability space, or background field
is assumed.
\end{standingrule}

\begin{standingrule}[\status{D}]\label{sr:licensing}
A structure may appear in a theorem only when its generating conditions
have been derived without presupposition.
\end{standingrule}

\begin{standingrule}[\status{D}]\label{sr:no-smuggling}
No observable distinction may be drawn by appealing to external naming,
primitive coordinates, or hidden memory of raw labels.
\end{standingrule}

\begin{remark}[\status{R}]
Standing Rules~\ref{sr:no-primitive}--\ref{sr:no-smuggling} are not
denials that geometric, continuum, or probabilistic language may later
become lawful.  They are the claim that none of it is primitive at the
base of the spine.  Such language must be \emph{earned} through
certified transfer; it cannot be smuggled into the primitive grammar.
\end{remark}

%% -- Collapse Gates --------------------------------------------------

The following six conditions are the \emph{framework-collapse gates}:
structural pre-conditions whose violation makes admissible
observational comparison impossible.  Any finite relational system
failing one of these gates collapses: no stable quotient, no
well-defined defect ledger, and no downstream canonical force.

\begin{standingrule}[\status{D}]\label{sr:cg-alphabet}
\textup{(CG-1: Finite regime alphabet.)}
The set of stabilized collar-type classes is finite.  An infinite
collar-type alphabet prevents stabilization of the observational
quotient and collapses the framework.
\end{standingrule}

\begin{standingrule}[\status{D}]\label{sr:cg-fibers}
\textup{(CG-2: Finite fibers.)}
Every equivalence class in the observational quotient is finite.
Infinite fibers destroy the defect arithmetic and collapse the
bookkeeping.
\end{standingrule}

\begin{standingrule}[\status{D}]\label{sr:cg-projection}
\textup{(CG-3: Terminating projection.)}
The quotient projection $\pi: \Omega \to \Sigma_n$ has no
non-terminating reduction chain.  Non-termination collapses
the stabilization argument.
\end{standingrule}

\begin{standingrule}[\status{D}]\label{sr:cg-no-id-smuggle}
\textup{(CG-4: No identity smuggling.)}
No admissible test may produce a canonical inverse to the quotient
map at a singular interface.  Identity smuggling at singular points
re-introduces presentation-level surplus into the canonical grammar.
\end{standingrule}

\begin{standingrule}[\status{D}]\label{sr:cg-nesting}
\textup{(CG-5: Finite nesting depth.)}
The fibrational nesting depth of any licensed configuration is
bounded.  Infinite nesting prevents the finite stabilization argument
and collapses the spine.
\end{standingrule}

\begin{standingrule}[\status{D}]\label{sr:cg-valence}
\textup{(CG-6: Finite valence.)}
Every site in every licensed configuration has finitely many
relational neighbours.  Infinite valence collapses the local
observational structure on which the defect bookkeeping depends.
\end{standingrule}

\begin{remark}[\status{R}]
Standing Rules~\ref{sr:cg-alphabet}--\ref{sr:cg-valence} are not
independently postulated constraints; they follow from the requirement
that the framework is persistence-selectable.  Any system violating one
of these gates cannot sustain the admissible quotient construction and
exits the canonical program.  They are recorded here at the primitive
level so that later books may invoke them by name without returning to
the finiteness arguments.
\end{remark}

%% ---------------------------------------------------------------
\section{Finite Relational Arena}
%% ---------------------------------------------------------------

\begin{definition}[\status{D}]\label{def:signature}
A \emph{finite relational signature} is a finite family
\[
  L = \{R_i\}_{i \in I}
\]
in which each $R_i$ has finite arity $k_i \in \mathbb{N}$.
\end{definition}

\begin{definition}[\status{D}]\label{def:config}
A \emph{finite configuration} is a finite $L$-structure
\[
  \omega = \bigl(X,\, \{R_i^{\,\omega}\}_{i \in I}\bigr),
\]
where $X$ is a finite set of \emph{sites} and
$R_i^{\,\omega} \subseteq X^{k_i}$ for each $i \in I$.
\end{definition}

\begin{definition}[\status{D}]\label{def:config-space}
A \emph{configuration space} is a finite set $\Omega$ of finite
configurations over the same base signature $L$.
\end{definition}

\begin{definition}[\status{D}]\label{def:struct-equiv}
Two configurations $\omega, \omega' \in \Omega$ are
\emph{structurally equivalent}, written $\omega \cong \omega'$, if
there exists an $L$-isomorphism between them.
\end{definition}

\begin{remark}[\status{R}]\label{rmk:struct-vs-obs}
Structural equivalence compares raw relational presentations.  It does
not yet determine what admissible testing can distinguish.  Book~I is
not interested in presentation richness by itself; it is interested in
what survives admissible observability.  The objects that carry
canonical force are observational equivalence classes, not raw
isomorphism types.
\end{remark}

%% ---------------------------------------------------------------
\section{Admissible Processes, Tests, and Observational Equivalence}
%% ---------------------------------------------------------------

\begin{definition}[\status{D}]\label{def:ext-process}
An \emph{extension process} is a map
\[
  E: \Omega \to \mathrm{Str}^{\mathrm{fin}}_{L_E},
\]
where $L_E \supseteq L$ is a finite expansion of the base signature.
\end{definition}

\begin{definition}[\status{D}]\label{def:admissible-class}
A class $\mathcal{E}$ of extension processes is \emph{admissible} if
it satisfies the following four conditions.
\begin{enumerate}[label=\textup{(A\arabic*)}]
  \item\label{ax:A1} \emph{Closure under typed composition.}
        Whenever the composition of two processes in $\mathcal{E}$ is
        defined, the composite is also in $\mathcal{E}$.
  \item\label{ax:A2} \emph{Structural invariance.}
        $\omega \cong \omega'$ implies $E(\omega) \cong E(\omega')$ for
        all $E \in \mathcal{E}$.
  \item\label{ax:A3} \emph{Base-signature accountability.}
        Observational distinctions induced by $E$ remain readable at
        the base-signature level and do not smuggle external labels.
  \item\label{ax:A4} \emph{Finite output control.}
        For every $E \in \mathcal{E}$, the image $E(\Omega)$ has
        finitely many isomorphism types.
\end{enumerate}
\end{definition}

\begin{remark}[\status{R}]
Axiom~\ref{ax:A3} is the anti-smuggling gate at the process level.
A process that produces external labels or hidden coordinates violates
this axiom and is not admissible.  The formal identity extension may
serve as a composition unit, but it is not automatically a licensed
observational test; permitting it as such would allow primitive
re-identification at the ground floor.
\end{remark}

\begin{remark}[\status{R}]
\textup{(Architectural note: v6 condition~(iii).)}
An earlier form of this framework included a condition ``No New
Invariant Partitions'': any partition invariant under all
extensions must be determined by the base incidence structure.
In the canonical formulation this is split into
Theorem~\ref{thm:emax-unique} (maximality as theorem, not axiom)
and Axiom~(A3) (anti-smuggling content).
\end{remark}

\begin{remark}[\status{R}]
\textup{(Architectural note: v6 condition~(iii).)}
An earlier form of this framework included a condition ``No New
Invariant Partitions'': any partition of $\Omega$ invariant under all
extensions in $\mathcal{E}$ must be determined by the base incidence
structure.  In the canonical formulation this condition is split:
its maximality content is captured by Theorem~\ref{thm:emax-unique}
(a theorem, not an axiom), and its anti-smuggling content by
Axiom~(A3).  The canonical axiom set A1--A4 is therefore strictly
cleaner.
\end{remark}

\begin{definition}[\status{D}]\label{def:test-family}
A \emph{licensed internal test family} at stage $n$ is a finite family
$\mathcal{T}_n$ of admissible relational processes whose outputs are
evaluated only through structures already licensed at stage $n$.
Tests in $\mathcal{T}_n$ may not appeal to external naming, primitive
coordinates, or hidden memory of raw labels.
\end{definition}

\begin{definition}[\status{D}]\label{def:obs-equiv}
For $\omega, \omega' \in \Omega$, write $\omega \sim_n \omega'$ if for
every test $T \in \mathcal{T}_n$, the outputs $T(\omega)$ and
$T(\omega')$ yield the same currently licensed observable outcome, or
are both undefined.  We call $\sim_n$ \emph{observational equivalence
at stage~$n$}.
\end{definition}

\begin{lemma}[\status{U}]\label{lem:obs-equiv-equiv}
Observational equivalence $\sim_n$ is an equivalence relation on
$\Omega$.
\end{lemma}

\begin{proof}
\textit{Reflexivity.}  Every test $T \in \mathcal{T}_n$ returns the
same outcome on $\omega$ as on itself.

\textit{Symmetry.}  The condition is symmetric in $\omega$ and
$\omega'$ by definition.

\textit{Transitivity.}  If no test in $\mathcal{T}_n$ separates
$\omega$ from $\omega'$, and no test separates $\omega'$ from
$\omega''$, then no test separates $\omega$ from $\omega''$.
\end{proof}

\begin{definition}[\status{D}]\label{def:obs-quotient}
The \emph{observational quotient at stage~$n$} is
\[
  \Sigma_n \;:=\; \Omega/{\sim_n}.
\]
\end{definition}

\begin{definition}[\status{D}]\label{def:proj-interface}
A \emph{projection interface} is a partial reduction map
$\pi: \Omega \dashrightarrow \Omega_{\mathrm{red}}$ that preserves
currently licensed observable behavior while forgetting finer
relational structure.
\end{definition}

\begin{definition}[\status{D}]\label{def:singular-interface}
A projection interface $\pi$ is \emph{singular at $y \in
\Omega_{\mathrm{red}}$} if there exist distinct $\omega_1, \omega_2
\in \Omega$ with $\pi(\omega_1) = \pi(\omega_2) = y$, and no licensed
inverse selection rule at $y$ distinguishes $\omega_1$ from $\omega_2$.
\end{definition}

\begin{proposition}[\status{U}]\label{prop:singular}
If there exist $\omega_1 \neq \omega_2$ in $\Omega$ with $\omega_1
\sim_n \omega_2$, then the canonical quotient map
$q: \Omega \to \Sigma_n$ is singular at the class
$[\omega_1]_{\sim_n} = [\omega_2]_{\sim_n}$.
\end{proposition}

\begin{proof}
Because $\omega_1 \sim_n \omega_2$, both map to the same quotient
class.  Because $\omega_1 \neq \omega_2$, the map is many-to-one at
that class.  Any admissible inverse selector would distinguish
$\omega_1$ from $\omega_2$, contradicting observational equivalence.
\end{proof}

\begin{remark}[\status{R}]\label{rmk:singular-force}
Proposition~\ref{prop:singular} records an inescapable structural
feature: wherever the observational quotient map is many-to-one,
canonical inverse selection is impossible without illicit
discrimination.  The absence of a canonical inverse is not a defect
of formalization; it is a genuine structural fact that later books
must respect.  All certified downstream structure must factor through
the quotient, not through presentation-level surplus.
\end{remark}

\begin{corollary}[\status{U}]\label{cor:quotient-welldef}
The observational quotient $\Sigma_n$ is well-defined, and canonical
observational content at stage~$n$ is determined by admissible test
equivalence, not by raw token identity or structural equivalence.
\end{corollary}

\begin{proof}
Lemma~\ref{lem:obs-equiv-equiv} ensures $\sim_n$ is an equivalence
relation, so $\Omega/{\sim_n}$ is well-defined.  The second clause
follows from the definition of $\sim_n$: the only distinctions that
survive to $\Sigma_n$ are those that some licensed test in
$\mathcal{T}_n$ can detect.
\end{proof}

%% -- Behavioral Congruence Canonicity --------------------------------

\begin{lemma}[\status{U}]\label{lem:behavioral-canonical}
\textup{(Behavioral Congruence Canonicity.)}
Among all equivalence relations on $\Omega$ induced by test families
$\mathcal{T}'$ satisfying the admissibility constraints~\textup{(A1)--(A4)},
the observational equivalence $\sim_n$ is the finest: if $\omega \sim_n
\omega'$ then $\omega \approx \omega'$ for every such $\approx$.
\end{lemma}

\begin{proof}
Any test family $\mathcal{T}'$ satisfying (A1)--(A4) is a subset of the
full admissible family $\mathcal{T}_n$.  If $\omega$ and $\omega'$ are
indistinguishable by every test in $\mathcal{T}_n$, they are in
particular indistinguishable by every test in $\mathcal{T}' \subseteq
\mathcal{T}_n$.  Hence $\omega \sim_n \omega'$ implies $\omega \approx
\omega'$ for every admissibility-consistent~$\approx$.
\end{proof}

\begin{remark}[\status{R}]
Lemma~\ref{lem:behavioral-canonical} confirms that the observational
quotient $\Sigma_n$ is not an arbitrary choice: it is the canonical
finest partition consistent with the admissibility axioms.  No
admissibility-consistent equivalence relation can distinguish more
than $\sim_n$ distinguishes.  This is the operational force of the
no-smuggling discipline at the level of equivalence structure.
\end{remark}

%% -- Admissible-Class Canonicity -------------------------------------

\begin{theorem}[\status{U}]\label{thm:emax-unique}
\textup{(Admissible-Class Canonicity.)}
Let $(\Omega, L)$ be a finite relational substrate with finite
compression space $|\Sigma_\infty| < \infty$.  Define
\begin{equation*}
\begin{split}
  E_{\max} \;:=\; \bigl\{\, E : \Omega \to \mathrm{Str}^{\mathrm{fin}}
  \;\big|\;\\
  &\quad E \text{ satisfies (A1)--(A4) and }
  E(\Omega) \text{ has fin.~many isom.~types}\,\bigr\}.
\end{split}
\end{equation*}
Then $E_{\max}$ is the unique maximal admissible class: every admissible
class $\mathcal{E}$ satisfies $\mathcal{E} \subseteq E_{\max}$.
\end{theorem}

\begin{proof}
By definition, $E_{\max}$ consists of every map satisfying (A1)--(A4)
together with finite output control.  Any admissible class
$\mathcal{E}$ satisfies (A1)--(A4) by definition, so $\mathcal{E}
\subseteq E_{\max}$.  Uniqueness is immediate since $E_{\max}$ is
defined as the set of all such maps; any two descriptions of the
maximal admissible class must coincide.
\end{proof}

\begin{remark}[\status{R}]
Theorem~\ref{thm:emax-unique} shows that the admissibility axioms
determine a canonical ceiling on what counts as an admissible
extension.  Every concrete admissible class studied in the program is
a sub-class of $E_{\max}$.  No admissible test can exceed the bounds
set by (A1)--(A4) without exiting the canonical framework.
\end{remark}

%% -- K_0 Primality ---------------------------------------------------

\begin{standingrule}[\status{D}]\label{sr:alphabet-prime}
\textup{(Collar alphabet size $K_0$.)}
Let $K_0 \in \mathbb{N}$ denote the number of distinct stabilized
WL-1 collar-type classes of the canonical finite relational
substrate under the full admissible class $E_{\max}$.  This number is
a structural invariant of the substrate, determined by the
stabilized quotient $\Sigma_\infty$.
\end{standingrule}

\begin{theorem}[\status{U}]\label{thm:K0-prime}
\textup{(Primality and Uniqueness of the Collar Alphabet.)}
The collar alphabet size $K_0$ is prime.  Moreover, $K_0 = 7$ is the
unique prime satisfying the structural self-consistency condition
$F(p) = p$, where $F(p) := |\Phi(\mathfrak{g}_{\mathrm{color}}(p))| + 1$.
\end{theorem}

\begin{proof}
\textit{Primality.}  If $K_0 = n$ is composite then
$|\mathrm{Aut}(\mathbb{Z}_n)| = \varphi(n) < n-1$, so proper
non-trivial orbits exist; Axiom~(A2) forces their identification,
collapsing the alphabet below~$n$.  For prime $p$ all $p$ classes
are mutually distinguishable.  Hence $K_0$ is prime.

\textit{Uniqueness at $K_0 = 7$.}  Exhaustion over odd primes via the
pendant-algebra Chevalley isomorphism:
$p=3$: $F(3)=3=p$ but $\tau=3$, non-viable;
$p=5$: $F(5)=7\neq 5$;
$p=7$: $F(7)=7=p$, gauge algebra
$\mathfrak{su}(3)\oplus\mathfrak{su}(2)\oplus\mathfrak{u}(1)$,
$r=4=\tau$, all conditions satisfied;
$p\geq 11$: $r(p)>\tau(p)$.
(Discharged unconditionally in ledger revision~IR8, March~2026.)
\end{proof}

\begin{remark}[\status{R}]
The primality clause follows from Axiom~(A2) alone.  The uniqueness
clause ($K_0=7$) uses the Weyl structure identity and the fixed-point
condition $F(p)=p$; it requires no branch-specific hypotheses.  This
is a strictly stronger statement than the earlier form.
\end{remark}

%% ---------------------------------------------------------------
\section{Refinement and Stabilization}
%% ---------------------------------------------------------------

Let $\mathcal{T}_0 \subseteq \mathcal{T}_1 \subseteq \mathcal{T}_2
\subseteq \cdots$ be a nested sequence of licensed test families.
Each stage induces an observational equivalence $\sim_n$ and quotient
$\Sigma_n = \Omega/{\sim_n}$.

\begin{definition}[\status{D}]\label{def:refinement-chain}
A \emph{refinement chain} on $\Omega$ is a nested sequence of licensed
test families $(\mathcal{T}_n)_{n \geq 0}$ such that each induced
quotient $\Sigma_n$ is finite.
\end{definition}

\begin{proposition}[\status{U}]\label{prop:monotone}
\textup{(Monotonicity.)}  If $\mathcal{T}_n \subseteq
\mathcal{T}_{n+1}$, then $\omega \sim_{n+1} \omega'$ implies $\omega
\sim_n \omega'$.  Equivalently, every class of $\Sigma_{n+1}$ is
contained in a class of $\Sigma_n$.
\end{proposition}

\begin{proof}
If two configurations are observationally indistinguishable by the
larger test family $\mathcal{T}_{n+1}$, they are indistinguishable by
every test in the smaller family $\mathcal{T}_n$.  Hence every
equivalence class for $\sim_{n+1}$ lies inside an equivalence class
for $\sim_n$.
\end{proof}

\begin{remark}[\status{R}]
Proposition~\ref{prop:monotone} means refinement proceeds only by
partition splitting: classes may stay the same or split into smaller
classes, but they never merge.  This monotone splitting is the
structural engine behind finite stabilization.
\end{remark}

\begin{definition}[\status{D}]\label{def:stabilization}
A refinement chain \emph{stabilizes at stage $N$} if
$\Sigma_n = \Sigma_N$ for all $n \geq N$.
\end{definition}

\begin{theorem}[\status{U}]\label{thm:stabilization}
\textup{(Finite Stabilization Theorem.)}  Every refinement chain on a
finite configuration space $\Omega$ stabilizes after finitely many
steps.
\end{theorem}

\begin{proof}
Each quotient $\Sigma_n$ corresponds to a partition of the finite set
$\Omega$.  By Proposition~\ref{prop:monotone}, these partitions are
successively finer: no two classes merge.  The total number of
distinct partitions of a finite set is finite (bounded by the Bell
number $B_{|\Omega|}$).  Therefore a strictly refining sequence of
partitions cannot continue indefinitely: once every possible splitting
has occurred, the sequence must stabilize.  Hence there exists a stage
$N$ such that $\Sigma_n = \Sigma_N$ for all $n \geq N$.
\end{proof}

\begin{definition}[\status{D}]\label{def:stabilized-quotient}
Fix a stabilization stage $N$ and define the
\emph{stabilized observational quotient}
\[
  \Sigma_\infty \;:=\; \Sigma_N.
\]
\end{definition}

\begin{remark}[\status{R}]\label{rmk:sigma-infty}
The notation $\Sigma_\infty$ is mnemonic, not ontological.  It records
that further admissible refinement no longer changes the quotient.  It
does not assert that a primitive infinite-limit object has been
introduced; no such object exists at this stage.  The subscript
$\infty$ means \emph{eventual stability}, nothing more.
\end{remark}

\begin{corollary}[\status{U}]\label{cor:stabilized-welldef}
The stabilized observational quotient $\Sigma_\infty$ is well-defined
and independent of the particular stabilization stage $N$ chosen.
\end{corollary}

\begin{proof}
By Theorem~\ref{thm:stabilization}, a stabilization stage $N$ exists.
If $N'$ is another stabilization stage with $N' \geq N$, then
$\Sigma_{N'} = \Sigma_N$ by the definition of stabilization.
\end{proof}

%% ---------------------------------------------------------------
\section{Survivor Chains and Defect Bookkeeping}
%% ---------------------------------------------------------------

Before Book~I ends, it introduces the primitive bookkeeping that
later supports Book~III's transport language.  At this stage, these
are purely combinatorial objects with no dynamical or physical
interpretation.

\begin{definition}[\status{D}]\label{def:survivor-component}
Let $C_n \in \Sigma_n$.  A \emph{survivor chain} issuing from $C_n$
is a nested sequence of classes
\[
  C_n \supseteq C_{n+1} \supseteq C_{n+2} \supseteq \cdots,
  \quad C_m \in \Sigma_m \text{ for all } m \geq n.
\]
A survivor chain records class-level persistence through successive
admissible refinements.  It is not a token-level identity claim; it is
a quotient-class-level persistence claim compatible with compression.
\end{definition}

\begin{definition}[\status{D}]\label{def:continuation-defect}
Let $C_n \in \Sigma_n$.  The \emph{continuation defect} of $C_n$ at
stage~$n$ is the failure of stage-$n$ admissible data to determine a
unique descendant class in $\Sigma_{n+1}$.
\end{definition}

\begin{definition}[\status{D}]\label{def:one-step-defect}
For $C_n \in \Sigma_n$, define the \emph{one-step defect number}
\[
  \delta(C_n) \;:=\; \#\{C_{n+1} \in \Sigma_{n+1} : C_{n+1} \subseteq C_n\}
  \;-\; 1.
\]
Thus $\delta(C_n) = 0$ exactly when continuation to the next stage is
unique.
\end{definition}

\begin{proposition}[\status{U}]\label{prop:defect-nonneg}
For every $C_n \in \Sigma_n$, one has $\delta(C_n) \geq 0$, with
equality if and only if $C_n$ has a unique descendant in $\Sigma_{n+1}$.
\end{proposition}

\begin{proof}
Every element of $C_n$ belongs to some class of the finer partition
$\Sigma_{n+1}$, so at least one descendant exists.  Therefore the
number of descendants is at least~$1$, hence $\delta(C_n) \geq 0$.
Equality holds exactly when there is precisely one descendant.
\end{proof}

\begin{definition}[\status{D}]\label{def:defect-ledger}
For a survivor chain $C_n \supseteq C_{n+1} \supseteq \cdots
\supseteq C_m$, its \emph{defect ledger} is
\[
  \Delta(C_n \leadsto C_m) \;:=\; \sum_{k=n}^{m-1} \delta(C_k).
\]
This records the cumulative branching burden along the chain.
\end{definition}

\begin{proposition}[\status{U}]\label{prop:ledger-additivity}
\textup{(Ledger additivity.)}  For any intermediate stage $\ell$ with
$n \leq \ell \leq m$,
\[
  \Delta(C_n \leadsto C_m)
  \;=\;
  \Delta(C_n \leadsto C_\ell) + \Delta(C_\ell \leadsto C_m).
\]
\end{proposition}

\begin{proof}
Immediate from the definition of $\Delta$ as a finite sum of one-step
defect numbers: splitting the sum at index $\ell-1$ gives the two
partial sums on the right.
\end{proof}

%% -- Structural Entropy ----------------------------------------------

\begin{definition}[\status{D}]\label{def:structural-entropy}
The \emph{structural entropy} of a configuration space $\Omega$ (with
respect to the stabilized observational quotient $\Sigma_\infty$) is
\[
  S(\Omega) \;:=\; \log|\Sigma_\infty|.
\]
No thermodynamic, probabilistic, or physical interpretation is assumed.
$S(\Omega)$ measures the logarithmic number of admissible distinctions
the stabilized quotient can sustain: the \emph{distinguishability
capacity} of the system.
\end{definition}

\begin{theorem}[\status{U}]\label{thm:entropy-monotone}
\textup{(Entropy Monotonicity.)}
For any right-stable admissible extension $E : \Omega \to \Omega'$
(one that does not create new defect events),
\[
  S(E(\Omega)) \;\leq\; S(\Omega).
\]
For defect-mediated extensions (those that do create new defects), the
defect ledger $\Delta$ records the entropy budget consumed:
\[
  S(E(\Omega)) \;\leq\; S(\Omega) + \Delta(\text{chain through } E).
\]
\end{theorem}

\begin{proof}
A right-stable extension maps each equivalence class of $\Sigma_\infty$
into at most one equivalence class of the refined quotient, so the
number of classes cannot increase: $|\Sigma'_\infty| \leq |\Sigma_\infty|$,
hence $S(E(\Omega)) \leq S(\Omega)$.

For a defect-mediated extension, each defect event records exactly one
class-splitting event.  The cumulative defect $\Delta$ bounds the
total number of class splits; hence
$|\Sigma'_\infty| \leq |\Sigma_\infty| \cdot 2^{\Delta}$
and $S(E(\Omega)) \leq S(\Omega) + \Delta \log 2 \leq S(\Omega) + \Delta$.
\end{proof}

\begin{remark}[\status{R}]
Theorem~\ref{thm:entropy-monotone} has two faces: right-stable
extensions cannot increase entropy (the quotient only coarsens or
stays put), while defect-mediated extensions may increase entropy but
only within bounds recorded by the defect ledger.  The defect ledger
is therefore the complete bookkeeper of entropic budget in the program.
\end{remark}

\begin{theorem}[\status{U}]\label{thm:ledger-entropy-duality}
\textup{(Ledger--Entropy Duality.)}
For any survivor chain $C_0 \supseteq C_1 \supseteq \cdots
\supseteq C_n$,
\[
  S(C_0) - S(C_n) \;=\; \Delta(C_0 \leadsto C_n).
\]
\end{theorem}

\begin{proof}
Each step: $S(C_k)-S(C_{k+1})=\delta(C_k)$ by
Theorem~\ref{thm:entropy-monotone}.  Summing and applying
Proposition~\ref{prop:ledger-additivity} gives the result.
\end{proof}

\begin{remark}[\status{R}]
Definition~\ref{def:structural-entropy} (running balance) and
Definition~\ref{def:defect-ledger} (transaction log) are two
readings of the same object.
\end{remark}

%% -- Irreversible Temporal Preorder (unconditional arm only) ---------

\begin{scholium}[\status{U}]\label{sch:temporal-preorder}
\textup{(Irreversible Preorder from Defect Ledger.)}
Define a preorder on the history space $\mathrm{Hist}(R)$ by
\[
  h \preceq h' \;\iff\; \Delta(h) \leq \Delta(h'),
\]
where $\Delta(h)$ is the cumulative defect ledger weight along history
$h$.  Then $(\mathrm{Hist}(R), \preceq)$ is a preorder:
\begin{enumerate}[label=\textup{(\roman*)}]
  \item \emph{Reflexivity and transitivity.}  Immediate from $\leq$ on
        $\mathbb{N}$.
  \item \emph{Forward-only.}  $\Delta$ is non-decreasing along any
        admissible continuation, since each transition appends a
        non-negative one-step defect $\delta(C_n) \geq 0$
        (Proposition~\ref{prop:defect-nonneg}).  No admissible
        continuation can decrease $\Delta$.
\end{enumerate}
The quotient $\mathrm{Hist}(R)/{\equiv}$ (where $h \equiv h'$ iff
$\Delta(h) = \Delta(h')$) carries a partial order with antisymmetry by
definition.

\emph{This is the unconditional arm of the temporal preorder.}  The
conditional arm — interpreting $\preceq$ as a physical temporal order
requiring Phase-A stability — is a theorem of Book~II
(Book~II, Scholium~II.9.x).  No metric, rate,
simultaneity structure, or causal propagation speed is derived here.
\end{scholium}

%% -- Holographic Sparsity (conditional, brief) -----------------------

\begin{corollary}[\status{C}]\label{cor:holographic-sparsity}
\textup{(Holographic Sparsity.)}
\textup{[Conditional on WL-1 stabilization package and LS-2 (imported
from Book~II).]}
For a finite connected region $X$ with $m$ pairwise
$\Delta_{\mathrm{sep}}$-separated fork sites,
\[
  m \log 2 \;\leq\; \log W(X) \;\leq\; (\log K_0)\,|\partial X|,
\]
where $W(X)$ is the number of distinct witness outcomes on $X$ and
$\partial X$ is the boundary of $X$.  The upper bound is the interface
capacity inequality~IC-1 (Book~II); the lower bound is the fork
entropy corollary.  The area-law character is a structural consequence
of finite alphabet and LS-2, not an imported physical principle.
\end{corollary}

\begin{proof}
Upper bound: by IC-1, each boundary site can contribute at most
$\log K_0$ bits of distinct witnessable outcome; with $|\partial X|$
boundary sites the bound follows.  Lower bound: the $m$ fork sites are
pairwise separated, so each contributes an independent binary choice to
$W(X)$, yielding $W(X) \geq 2^m$.
\end{proof}

\begin{remark}[\status{R}]
Corollary~\ref{cor:holographic-sparsity} is conditional on LS-2, a
hypothesis of Book~II.  It is stated here for architectural completeness
because its upper bound (IC-1) is a direct consequence of the primitive
finite-alphabet constraint $K_0 < \infty$ established in this book.
Full proof requires the collar machinery of Book~II.
\end{remark}

%% -- Computational Scholia: G_9 Barbell Analysis ---------------------

\begin{scholium}[\status{D}]\label{sch:ledger-quotient-G9}
\textup{(Canonical Ledger Quotient on $G_9$.)}  \textup{[Computational
certification.  Analytic proof open.]}
Let $G_9 = K_3\text{-}b_3\text{-}K_3$ be the barbell graph on 9
vertices with $V \cong \mathbb{F}_2^4$ as the WL-1 color space
(4 bits: midpoint, connectors, bridge ends, clique bulk).  Among the
28 radius-2 single-edge extensions from the base configuration:
\begin{enumerate}[label=\textup{(\arabic*)}]
  \item 22 preserve all 14 nontrivial orbit-constant refinements
        (defect-free).
  \item 6 break exactly 8 of the 14 refinements (bimodal obstruction).
  \item All 6 obstructing extensions share the same survivor subspace
        $W \subset V$ with $|W| = 8$.
  \item The unique linear constraint is
        $W = \{x \in V : \mathrm{bit}_2 + \mathrm{bit}_4 \equiv 0
        \pmod{2}\}$.
  \item The annihilator $a = (0,1,0,1) \in V^*$ is the unique
        codimension-1 kernel direction (the \emph{canonical ledger
        axis}).
\end{enumerate}
Source: v7.35 Theorem~55.1 (exhaustive BFS enumeration).
\end{scholium}

\begin{scholium}[\status{D}]\label{sch:depth2-ecology-G9}
\textup{(Depth-2 Context Ecology on $G_9$.)}  \textup{[Computational
certification.]}
Under BFS to depth~2 (495 states, 852 transitions) from $G_9$, the
reachable configurations partition into 8 distinct survivor contexts
stratified by dimension as $1 + 3 + 3 + 1$:
\begin{enumerate}[label=\textup{(\arabic*)}]
  \item Dimension~1: 1 context (10 states).
  \item Dimension~2: 3 contexts (16, 14, 10 states).
  \item Dimension~3: 3 contexts (109, 6, 2 states) — the
        \emph{three-family structure}.
  \item Dimension~4: 1 context (328 states, full $V$).
\end{enumerate}
The dimension-3 triple provides computational evidence that three
distinct survivor families coexist under the admissible dynamics at
this substrate scale.  Source: v7.35 Theorem~56.1.
\end{scholium}

\begin{scholium}[\status{D}]\label{sch:gating-depths-G9}
\textup{(Certified Gating Depths on $G_9$.)}  \textup{[Computational
certification.]}
On $G_9$ with $V \cong \mathbb{F}_2^4$, the six weight-2
pair-identification hyperplanes have certified gating depths:
\begin{center}
\begin{tabular}{cc}
\hline
Constraint & Gating depth \\
\hline
$\mathrm{bit}_2 + \mathrm{bit}_4 = 0$ & 1 \\
$\mathrm{bit}_3 + \mathrm{bit}_4 = 0$ & 2 \\
$\mathrm{bit}_1 + \mathrm{bit}_3 = 0$ & 2 \\
$\mathrm{bit}_1 + \mathrm{bit}_2 = 0$ & 2 \\
\hline
\end{tabular}
\end{center}
The depth-1 gate $(\mathrm{bit}_2 + \mathrm{bit}_4 = 0)$ is the
canonical ledger axis identified in
Scholium~\ref{sch:ledger-quotient-G9}.  Source: v7.35
Proposition~57.3.
\end{scholium}

\begin{definition}[\status{D}]\label{def:stabilizer-class}
A class $C_n \in \Sigma_n$ is a \emph{stabilizer class} if every
survivor chain issuing from $C_n$ has zero defect from stage~$n$
onward.
\end{definition}

\begin{proposition}[\status{U}]\label{prop:stabilizer-after-stab}
If the refinement chain stabilizes globally at stage $N$, then every
class in $\Sigma_N$ is a stabilizer class.
\end{proposition}

\begin{proof}
After stage $N$, the partitions do not change, so every class has
exactly one descendant at every later stage: itself.  Hence every
later one-step defect vanishes.
\end{proof}

\begin{remark}[\status{R}]\label{rmk:defect-pre-transport}
At the level of Book~I, the defect ledger is purely combinatorial
bookkeeping.  The concepts of continuation, persistence across
contexts, transport obstruction, and finite balance will not acquire
their full force until Book~III.  The purpose of introducing the
defect ledger here is to fix its arithmetic substrate cleanly at the
primitive level, before any dynamical or comparative structure is
available.  Later books may invoke these objects through their
canonical names without returning to the primitive setup.
\end{remark}

%% ---------------------------------------------------------------
\section{Governing Theorem of Book I}
%% ---------------------------------------------------------------

The content of Book~I can now be compressed into its governing theorem.

\begin{theorem}[\status{U}]\label{thm:governing}
\textup{(Observational Quotient Theorem.)}  Finite admissible-test
quotients define a coherent observational state space, and every
finite refinement chain stabilizes after finitely many steps.

More precisely:
\begin{enumerate}[label=\textup{(\alph*)}]
  \item\label{thm:gov-a} The observational quotient $\Sigma_n$ is
        well-defined at every stage and canonical observational content
        is determined by admissible test equivalence, not by raw token
        identity.
  \item\label{thm:gov-b} Every refinement chain on a finite
        configuration space stabilizes: there exists $N$ such that
        $\Sigma_n = \Sigma_N =: \Sigma_\infty$ for all $n \geq N$.
\end{enumerate}
\end{theorem}

\begin{proof}
Part~\ref{thm:gov-a} is Corollary~\ref{cor:quotient-welldef}.
Part~\ref{thm:gov-b} is Theorem~\ref{thm:stabilization}.
\end{proof}

\begin{corollary}[\status{U}]\label{cor:no-presentation-surplus}
No later canonical theorem in the spine may depend on
presentation-level distinctions that are invisible to the stabilized
observational quotient.
\end{corollary}

\begin{proof}
By Theorem~\ref{thm:governing}\ref{thm:gov-a}, the only distinctions
with canonical force are those detectable by admissible tests.
Presentation-level distinctions collapsed by $\sim_\infty$ carry no
admissible witness, and therefore no canonical force.
\end{proof}

\begin{corollary}[\status{U}]\label{cor:defect-welldef}
The defect ledger is a well-defined combinatorial invariant of survivor
chains, additive along concatenation, and vanishes on every survivor
chain after the global stabilization stage.
\end{corollary}

\begin{proof}
Propositions~\ref{prop:defect-nonneg}, \ref{prop:ledger-additivity},
and~\ref{prop:stabilizer-after-stab}.
\end{proof}

\begin{scholium}[\status{R}]\label{sch:book-i-boundary}
Book~I establishes only the primitive theorem-bearing kernel: admissible
tests, observational quotient, finite stabilization, and defect
bookkeeping.  It does not prove local stability of observable outcomes,
boundary determinacy, source universality, branch legitimacy, or
continuum correspondence.  Every one of those claims belongs to a
later book and may not be silently imported here.
\end{scholium}

%% ---------------------------------------------------------------
\section{Boundary of Book I}
%% ---------------------------------------------------------------

The following is an explicit statement of what Book~I has and has not
proved.

\textbf{What Book~I proves.}
\begin{enumerate}
  \item Every finite configuration space admits a coherent
        observational quotient induced by any licensed test family.
  \item Observational equivalence is an equivalence relation; the
        quotient map is many-to-one precisely where distinct
        configurations are observationally indistinguishable, forcing
        singular projection interfaces.
  \item Every finite refinement chain stabilizes after finitely many
        steps.  The stabilized quotient $\Sigma_\infty$ is well-defined.
  \item The defect ledger is a well-defined additive combinatorial
        invariant of survivor chains; it vanishes after global
        stabilization.
  \item The six collapse gates (CG-1 through CG-6) are the
        pre-conditions for any persistent admissible object.
  \item Behavioral congruence $\sim_n$ is the canonical finest
        admissibility-consistent equivalence; the maximal admissible
        class $E_{\max}$ is unique (Thm~\ref{thm:emax-unique}).
  \item The collar alphabet size $K_0$ is prime
        (Thm~\ref{thm:K0-prime}).
  \item Structural entropy $S(\Omega) = \log|\Sigma_\infty|$ is
        well-defined and monotone non-increasing for right-stable
        extensions (Thm~\ref{thm:entropy-monotone}).
  \item The defect ledger defines an irreversible preorder on
        continuation histories (Scholium~\ref{sch:temporal-preorder}).
\end{enumerate}

\textbf{What Book~I does not prove, and may not be assumed here.}
\begin{enumerate}
  \item \emph{Locality.}  No claim is made that observable outcomes
        in a region are controlled by boundary data.  That requires
        the LS-2 package of Book~II.
  \item \emph{Boundary capacity.}  No counting bounds on outcome
        multiplicity are derived here.
  \item \emph{Source universality.}  The existence of a universal
        source object is a theorem of Book~IV, not of Book~I.
  \item \emph{Persistence and transport.}  The transport and
        persistence machinery of Book~III does not exist here.  The
        defect ledger is primitive bookkeeping, not a transport law.
  \item \emph{Continuum structure.}  No continuum, differential,
        integral, or smooth language is licensed at this stage.
  \item \emph{Physical interpretation.}  No identification with
        spacetime, gauge fields, matter, or any physical theory is
        claimed or implied.
\end{enumerate}

%% ---------------------------------------------------------------
\section{Handoff to Book II}
%% ---------------------------------------------------------------

Book~II begins when the stabilized observational structure established
here is localized.  The natural next question — once one has a
stabilized observational quotient $\Sigma_\infty$ — is whether the
observable outcome on a finite subregion is determined by what happens
at its boundary.  That question requires introducing regions, collars,
and stabilized collar types: all direct descendants of the objects
defined in this book.

The transition from Book~I to Book~II is:
\begin{gather*}
\text{adm.~process class}
\to\text{obs.~quotient}
\to\text{finite stabilization}\\
\to\text{stabilized local types and boundary law.}
\end{gather*}
Book~I supplies the first three terms of this chain.
Book~II supplies the fourth.

%% ---------------------------------------------------------------
\section{Dependency Ledger}
%% ---------------------------------------------------------------

The following table records the load-bearing items of Book~I, their
status, and their export destinations.

\medskip
\renewcommand{\arraystretch}{1.4}
\noindent
\begin{tabular}{@{}p{3.6cm}p{0.8cm}p{3.8cm}p{4.8cm}@{}}
\hline
\textbf{Theorem} & \textbf{St.} & \textbf{Depends on} & \textbf{Exports to} \\
\hline
\ref{lem:obs-equiv-equiv}\ Obs.\ equiv.\ rel.
  & [U] & Def.~\ref{def:obs-equiv}
  & Defs.~\ref{def:obs-quotient}, \ref{thm:governing} \\

\ref{prop:singular}\ Singular interfaces
  & [U] & \ref{lem:obs-equiv-equiv}, \ref{def:singular-interface}
  & Rmk.~\ref{rmk:singular-force}; Book~V \\

\ref{cor:quotient-welldef}\ Quotient well-def.
  & [U] & \ref{lem:obs-equiv-equiv}
  & \ref{thm:governing}; Books~II--VII \\

\ref{prop:monotone}\ Monotonicity
  & [U] & \ref{def:refinement-chain}
  & \ref{thm:stabilization} \\

\ref{thm:stabilization}\ Finite Stabilization
  & [U] & \ref{prop:monotone}, fin.\ of $\Omega$
  & \ref{thm:governing}; Books~II--VII \\

\ref{cor:stabilized-welldef}\ $\Sigma_\infty$ well-def.
  & [U] & \ref{thm:stabilization}
  & Books~II--VII \\

\ref{prop:defect-nonneg}\ $\delta(C_n)\geq 0$
  & [U] & \ref{def:one-step-defect}
  & \ref{prop:ledger-additivity}, \ref{thm:governing} \\

\ref{prop:ledger-additivity}\ Ledger additivity
  & [U] & \ref{def:defect-ledger}
  & \ref{cor:defect-welldef}; Book~III \\

\ref{prop:stabilizer-after-stab}\ Stabilizer classes
  & [U] & \ref{thm:stabilization}, \ref{def:stabilizer-class}
  & \ref{cor:defect-welldef}; Book~III \\

\ref{thm:entropy-monotone}\ Entropy Monotonicity
  & [U] & \ref{def:structural-entropy}, \ref{prop:ledger-additivity}
  & \ref{sch:temporal-preorder}; Book~III \\

\ref{sch:temporal-preorder}\ Irrev.\ preorder
  & [U] & \ref{prop:defect-nonneg}, \ref{prop:ledger-additivity}
  & Book~II Scholium (Phase-A arm); Book~III \\

\ref{cor:holographic-sparsity}\ Holographic Sparsity
  & [C] & \ref{sr:alphabet-prime}; Book~II IC-1
  & Book~I boundary; GR Branch Book \\

\ref{thm:governing}\ Obs.\ Quotient Theorem
  & [U] & \ref{cor:quotient-welldef}, \ref{thm:stabilization}
  & Books~II--VII; all branch books \\

\ref{cor:no-presentation-surplus}\ No pres.\ surplus
  & [U] & \ref{thm:governing}
  & Books~II--VII; Book~VII \\

\ref{cor:defect-welldef}\ Defect ledger
  & [U] & \ref{prop:defect-nonneg}--\ref{prop:stabilizer-after-stab}
  & Book~III transport machinery \\

\ref{lem:behavioral-canonical}\ Behav.\ Congruence Canon.
  & [U] & \ref{cor:quotient-welldef}; A1--A4
  & \ref{thm:emax-unique}; Books~II--VII \\

\ref{thm:emax-unique}\ Admiss.-Class Canon.
  & [U] & \ref{lem:behavioral-canonical}; A1--A4
  & Books~II--VII; Book~IV \\

\ref{thm:K0-prime}\ $K_0$ primality
  & [U] & A2; \ref{sr:alphabet-prime}
  & Book~II collar machinery; YM Branch Book \\
\hline
\end{tabular}
\renewcommand{\arraystretch}{1}

\medskip

\begin{remark}[\status{R}]
Every unconditional result in this ledger depends only on the
standing definitions and the finiteness of $\Omega$.  No hypothesis
beyond the finite relational framework is invoked.  The one
conditional result, Corollary~\ref{cor:holographic-sparsity}
(Holographic Sparsity), imports LS-2 from Book~II; it is recorded here
for architectural completeness only.  This unconditionality is the
source of Book~I's load-bearing character: everything downstream in
the spine rests on a foundation that requires no further assumptions
to stand.
\end{remark}

\end{document}
